\documentclass[11pt]{article}
\usepackage[margin=2.7cm]{geometry}
\usepackage{amsmath,amssymb}
\usepackage{booktabs}
\usepackage{hyperref}
\usepackage{parskip}
\usepackage{xcolor}
\usepackage{listings}

\lstset{
  basicstyle=\ttfamily\small,
  columns=fullflexible,
  breaklines=true,
  frame=single,
  language=R
}

\newcommand{\code}[1]{\texttt{#1}}
\newcommand{\file}[1]{\textcolor{black!60}{\texttt{#1}}}

\title{PI3K-mini: A Positive and a Negative Result for FIM Null-Space
Regularization\\[2pt]
\large What a small, fully-controlled toy model taught us, and what it
predicts for the real HGF model}
\author{}
\date{Companion to \file{scripts/example\_toy\_pi3kmini\_fim\_fit.R},
\file{scripts/example\_toy\_pi3kmini\_fim\_plots.R}, and
\file{scripts/example\_toy\_pi3kmini\_fim\_eigenvalue\_report.R}}

\begin{document}
\maketitle
\vspace{-1.5em}

\section{Purpose}

This note records the reasoning, experiments, and conclusions of an extended
investigation into whether the \code{groupL2()} / \code{nullSpaceProjector()}
FIM null-space regularization (\file{fim-portugal/functions-portugal.R}) can
be made to work reliably, using a small, fully-transparent toy model built
specifically to test the method's own precondition under full control. The
model, ``PI3K-mini'', is a deliberate mechanistic miniature of the PI3K
structural non-identifiability already documented for the real HGF/HepG2
model in \file{HGF\_HepG2/04-Results/M-22.0-comparison/05-Writeup/
Technical-Notes\_PI3K-Symmetry-and-Regularization-Diagnostics.pdf}. The
investigation produced both a clean positive result (Section~\ref{sec:clean})
and, on pushing the same model toward more realistic data, a clean negative
one (Section~\ref{sec:sweep}) --- together with a precise, verified
mechanistic explanation for the transition between the two
(Sections~\ref{sec:eigendecomp}--\ref{sec:practical}), a proposed hybrid
regularizer that keeps the positive part without inheriting the negative part
(Section~\ref{sec:hybrid}), and an assessment of what this predicts for the
real, 96-parameter HGF model (Section~\ref{sec:real-model}).

\section{The model}

\subsection{Topology}

\begin{center}
\begin{tabular}{ll}
\toprule
Reaction & Role \\
\midrule
$A \to B$, rate $k_1 A$ & control block; directly observed; unrelated to the symmetry \\
$\varnothing \to U$, rate $k_{\text{gen}} A$ & generates hidden $U$ catalytically ($A$ not consumed) \\
$U \to \varnothing$, rate $k_{\text{degU}} U$ & $U$'s own decay; \emph{not} part of the symmetry \\
$M_p \to M_{pp}$, rate $k_{\text{use}} M_p U$ & only readout of $U$'s activity, via $\mathrm{obs}M_{pp}$ \\
\bottomrule
\end{tabular}
\end{center}

Three conditions (\code{cond1/cond2/cond3}) differ only in $U$'s initial
value ($U_{\text{cond1}}=1,\,U_{\text{cond2}}=3,\,U_{\text{cond3}}=6$),
mirroring \code{PI3K\_5/PI3K\_250/PI3K\_350} in the real model. $A_0=1$ and
$M_{p,0}=1$ are fixed (not fitted) to avoid two \emph{unrelated} accidental
scale symmetries ($A_0 \leftrightarrow \mathrm{scale}_A/\mathrm{scale}_B$ and
$M_{p,0} \leftrightarrow \mathrm{scale}_{Mpp}$) that this design would
otherwise introduce.

\subsection{The exact symmetry}

For any $c>0$:
\[
U_{\text{cond}i} \to c\,U_{\text{cond}i}\ (i=1,2,3), \qquad
k_{\text{gen}} \to c\,k_{\text{gen}}, \qquad
k_{\text{use}} \to k_{\text{use}}/c
\]
leaves $U(t)$ scaling exactly by $c$ (linear, homogeneous ODE in $U$) and
leaves $k_{\text{use}}U(t)$ --- the only combination $M_p/M_{pp}$ ever see ---
exactly invariant. Every observable trajectory is therefore exactly invariant
under this one-parameter family, verified numerically to solver precision
($\sim 10^{-7}$, Section~\ref{sec:clean}). $k_{\text{degU}}$ does \emph{not}
load on this symmetry (linear decay of $U$ needs no compensation), mirroring
why the real model's \code{k\_degrade\_PI3K} does not load on the real PI3K
symmetry either.

The 5-parameter group ($U_{\text{cond1}},U_{\text{cond2}},U_{\text{cond3}},
k_{\text{gen}},k_{\text{use}}$) spans a rank-1 null space (one non-identifiable
direction) and a rank-4 identifiable complement, spanned by the four
combinations $k_{\text{use}}U_{\text{cond1}}$, $k_{\text{use}}U_{\text{cond2}}$,
$k_{\text{use}}U_{\text{cond3}}$, $k_{\text{gen}}k_{\text{use}}$ --- the same
1-null-in-5 structure as the real PI3K group. This was a deliberate design
choice: the minimal 2-parameter AB-model ($k\cdot q$, \S5.2 of the thesis
draft) has \emph{no} identifiable complement at all (rank-1 null space inside
a 2-dimensional group exhausts the space), so it cannot show the distinction
between ``the non-identifiable direction'' and ``a raw parameter that happens
to participate in it'' that turns out to be central to this whole
investigation (Section~\ref{sec:practical}).

\subsection{An implementation pitfall}
\label{sec:pitfall}

\code{Y(..., attach.input = TRUE)} attaches \emph{all} raw states, including
the hidden $U$, to the prediction output. Building the simulated dataset from
that attached output without filtering down to the actual observables first
silently fits against $U$'s own (non-invariant) trajectory too, which
destroys the intended symmetry with no error or warning --- the Hessian's
eigenvalue spectrum simply looks like an ordinary sloppy spread, no gap
anywhere. \code{attach.input = FALSE} is required. (This cost real debugging
time: the first version of the fit showed no spectral gap at all, which
briefly looked like the toy model design itself was wrong before the true
cause was found by directly comparing predicted trajectories under a
hand-applied symmetry shift against what \code{normL2()} was actually being
handed as ``data''.)

\section{Result 1: a clean win, on dense/precise data}
\label{sec:clean}

With $\sigma=0.03$, 31 time points per condition (dense, precise data), 200
multistart fits per objective (\code{obj\_data}, \code{obj\_l2 =
obj\_data + constraintL2(\ldots)}, \code{obj\_fim = obj\_data +
groupL2(P\_null,\ldots)}), all three centered at the \emph{same} reference
point $\mu=-1$ in every log-parameter (a fair-comparison fix --- see
Section~\ref{sec:same-center}):

\begin{center}
\begin{tabular}{lrrr}
\toprule
 & no regularization & FIM (group) & L2 (uniform) \\
\midrule
raw data $\chi^2$ at best fit & $-1186.38$ & $-1186.38$ & $-1186.28$ \\
identified-block RMSE & $0.099$ & $0.099$ & $0.159$ \\
invariant-combo RMSE & $0.061$ & $0.061$ & $0.081$ \\
SD of raw $U_{\text{cond1}}$, near-optimal fits & $2.06$--$7.6$ & $\approx 10^{-4}$ & $\approx 10^{-4}$ \\
\bottomrule
\end{tabular}
\end{center}

FIM matches the unregularized fit's accuracy exactly (unbiased) while
matching L2's reproducibility exactly (bounded) --- strictly better than
either alternative on its own, in the regime where the method's own
precondition (one isolated exact symmetry, everything else genuinely sharp)
holds by construction. Eigenvalue spectrum (log10) at this data quality:
$5.26,\,5.16,\,4.51,\,3.90,\,3.58,\,2.94,\,2.68,\,2.04,\,0.56,\,-11.6$ --- one
clean $\sim\!12$-order-of-magnitude gap, nothing else close to it.

\section{Same regularization center for both methods}
\label{sec:same-center}

An earlier version anchored \code{groupL2()} at \code{best\_unreg} (the
single best unregularized fit). Along the exactly flat direction this anchor
is itself an arbitrary draw from the unregularized multistart's own spread
--- in one run it landed at $\log U_{\text{cond1}} = 2.84$ (true value $0$),
so all FIM-regularized fits collapsed to that arbitrary point instead of
anywhere near truth, while L2 (anchored at $\mu=0$, coincidentally close to
$U_{\text{cond1}}$'s true value $1$) looked ``better'' for a reason that had
nothing to do with either method's merit. Fix: anchor both penalties at the
\emph{same} $\mu_{\text{center}}=-1$ for every parameter. This is mathematically
required for a fair comparison and also exposed a useful fact: with a shared
center, L2 and FIM land at (numerically) the \emph{same} point along the null
direction, because L2's cost on the stiff/identified directions is constant
along that valley --- so minimizing L2 restricted to the valley reduces to
the identical 1-D problem FIM's rank-1 penalty solves directly.

Caveat worth remembering: the individual raw value a regularizer settles on
along a truly non-identifiable direction is not itself a meaningful quantity
--- it is set entirely by the chosen anchor, not by data. Comparing raw
per-parameter estimates between two differently-anchored regularizers is not
a fair test of anything.

\section{Two outlier local optima under FIM}
\label{sec:outliers}

At the dense-data setting, dMod's own \code{plotPars()} (step-detected via
\code{stepDetect()}) flags 2 of 200 FIM fits as genuinely worse: raw
(unpenalized) $\chi^2 = -1086.1$ and $-729.4$ vs. $-1186.4$ for the rest ---
$100$ and $457$ worse, both reported \code{converged=TRUE}, 29--38
iterations, not truncated runs. The FIM penalty contributes $\approx 0$ at
all three points, confirming these are genuine alternative local optima of
the raw likelihood, not penalty artifacts. Predicted curves
(\file{pi3kmini\_fim\_outliers.png}) show $\mathrm{obs}M_{pp}$ in
\code{cond2}/\code{cond3} saturating almost instantly instead of following
the data's gradual early rise --- a classic ``which rate governs the fast vs.
slow phase'' ambiguity in fitting a hidden-species-driven rise, unrelated to
the exact symmetry (parameter shifts from the best fit are \emph{not}
consistent with a pure $c$-rescaling: $U_{\text{cond3}}$ moves independently
of $U_{\text{cond1}}/U_{\text{cond2}}/k_{\text{gen}}$). Neither unregularized
nor L2 hit this basin in the same 200-start sample (0/200 each); most simply
attributed to a small ($\sim\!1\%$) basin and ordinary multistart sampling
luck, not confirmed to be FIM-specific.

\section{Result 2: sparser/noisier data reverses the result}
\label{sec:sweep}

Reducing to 14 time points and raising $\sigma$ (later swept from $0.05$ to
$1.7$ at fixed data density) reverses the earlier ranking:

\begin{center}
\begin{tabular}{lrrr}
\toprule
 & no regularization & FIM (group) & L2 (uniform) \\
\midrule
identified-block RMSE & $1.19$ & $1.19$ & $0.38$ \\
invariant-combo RMSE & $6.36$ & $5.89$ & $0.33$ \\
raw $\chi^2$ at best fit & $-191.6$ & $-191.6$ & $-190.9$ \\
\bottomrule
\end{tabular}
\end{center}

L2 is now clearly \emph{better} than FIM on exactly the quantities FIM was
supposed to protect. The eigenvalue spectrum at 14 points shows a new,
second gap that did not exist at 31 points: $\ldots,0.36,\,{-1.25},\,{-2.44},
\,{-13.3}$ (log10) --- the true symmetry gap is untouched (still
$\sim\!11$ orders), but ranks 8--9 are new, genuinely emergent near-degenerate
directions.

\section{Diagnosing the new directions}
\label{sec:eigendecomp}

The rank-8 eigenvector ($\lambda\approx 10^{-1.25}$) is \emph{not} the true
symmetry: it loads heavily on $k_{\text{gen}}$ ($0.84$) and
$U_{\text{cond1/2/3}}$ ($-0.28$ each) with essentially \emph{no}
$k_{\text{use}}$ component ($-0.003$) --- a $k_{\text{gen}}$-vs-$U_{\text{cond}i}$
trade-off distinct from the true symmetry, which requires $k_{\text{use}}$ to
move oppositely. Projecting the near-optimal fits onto this new direction
vs. onto the true null direction (both computed from the data's own $H$)
gives a definitive result:

\begin{center}
\begin{tabular}{lrr}
\toprule
projection onto & unreg SD & FIM SD \\
\midrule
true exact symmetry (rank 10) & $2.32$ & $1.5\times 10^{-13}$ \\
new emergent direction (rank 8) & $2.51$ & $2.19$ \\
\bottomrule
\end{tabular}
\end{center}

FIM controls the direction it was built from essentially perfectly, and
provides \emph{zero} protection for the new one --- confirmed across every
saved $\sigma$: the true-symmetry-projection SD stays at
$10^{-13}$--$10^{-15}$ at \emph{every} $\sigma$ tested, unconditionally
(Section~\ref{sec:noise-scaling} explains why this holds exactly, not just
approximately).

\subsection{Correction: uniform noise does not reshape the spectrum}

A uniform $\sigma$ applied identically to every data point gives $H = S'S/
\sigma^2$ --- an exact positive scalar multiple of the same matrix at every
$\sigma$. Eigenvectors of a scalar-multiplied matrix are unchanged; only
eigenvalues scale. Checked directly (\file{pi3kmini\_eigenvalue\_report.csv},
\file{pi3kmini\_eigenvalues\_vs\_sigma.pdf}): consecutive-rank log10-gaps are
identical to 2 decimal places across every saved $\sigma$ (0.05--1.7) except
the last (true-symmetry) one, which is pure numerical noise around the exact
zero. \emph{The ``second gap'' found in Section~\ref{sec:sweep} came from
reducing the number of time points (which changes $S$ itself), not from
raising $\sigma$ (which only rescales $H$ uniformly and provably cannot
reshape relative gaps).} $P_{\text{null}}$ itself was verified numerically
identical (max abs.\ difference $6\times 10^{-13}$) between $\sigma=0.05$ and
$\sigma=1.7$.

\section{Why noise-driven curvature loss still hurts FIM}
\label{sec:noise-scaling}

Even with $P_{\text{null}}$ and the gap structure fixed, raising a uniform
$\sigma$ still shrinks \emph{every} data-only curvature entry as $1/\sigma^2$
--- including the four identifiable combinations inside the same 5-parameter
group. \code{constraintL2} adds a curvature of exactly $2/\sigma_{L2}^2$ to
\emph{every} coordinate, independent of $\sigma$; as data curvature shrinks
below this fixed floor, L2 increasingly (and correctly, in the sense of
``it always would'') dominates every coordinate, including the ones that
were never part of the symmetry. \code{groupL2} adds curvature only along
$P_{\text{null}}$'s span --- exactly zero elsewhere, at every $\sigma$ ---
so the identified complement's precision degrades with $\sigma$ exactly as
it would under no regularization at all. This is not a flaw specific to the
implementation; it is the necessary consequence of a penalty that is
deliberately inert outside one, exactly-known direction.

\section{Does pinning the null direction give unique parameter values?}
\label{sec:practical}

A natural expectation: if the null direction is regularized, doesn't the
whole 5-parameter block become uniquely (hence tightly) estimable? The
answer requires separating two different claims.

\textbf{True:} once $P_{\text{null}}$ is regularized, the full block has a
mathematically unique minimizer for a fixed dataset, because the other four
directions have real (if possibly small) curvature --- there is no longer an
infinite flat family.

\textbf{Not automatically true:} that this unique minimizer is found
\emph{precisely} by any real optimizer, or corresponds to a
\emph{practically} tight estimate. When curvature on those other directions
is small, the objective is nearly flat over a wide region around the unique
minimum: a small gradient (near any optimizer's stopping tolerance)
corresponds to a \emph{large} parameter displacement, since displacement
$\approx$ gradient / curvature. Different multistart runs land at genuinely
different points satisfying the same loose tolerance; even the exact
minimizer, if found, would carry a wide, honest confidence interval on that
direction.

This is exactly the \textbf{structural vs.\ practical identifiability}
distinction of Raue et al.\ (2009, already cited in the thesis bib). FIM's
projector removes only \emph{structural} non-identifiability (an exact,
dataset-independent symmetry). It has no mechanism against \emph{practical}
non-identifiability (a technically-unique but data-quality-dependent flat
region) --- and that is precisely what grows in the other four directions as
$\sigma$ increases. No method that targets only a fixed, precomputed exact
symmetry can fix that, because it is not a symmetry, only weak information.

\subsection{This is not special to PI3K-mini's 5-parameter group}

The identical logic applies to the AB-model's $k,q$: its ``identified''
direction ($k\cdot q$) has data-driven curvature that would also shrink under
worse data, and $k,q$ would show the same growing, \emph{correlated} scatter
(moving together along the weakening product direction, while staying locked
along the still-perfectly-regularized $k/q$ direction) that
\code{ab\_fit.R}'s own dense/precise dataset simply never demonstrated. The
2-parameter AB-model cannot show the richer phenomenon of Section~\ref{sec:eigendecomp}
(a raw parameter mixing a protected direction with unprotected ones), because
with only 2 parameters and 1 null direction there is no room for a third,
independent direction to emerge --- but the core mechanism (a fixed penalty
vs.\ a shrinking data curvature) is identical in both models.

\subsection{Summary statement for the chapter}

With sparser and noisier data, the model gets more \emph{sloppy}
(Gutenkunst et al., 2007): every direction's curvature shrinks, and
previously-comfortable directions sink below whatever curvature is needed to
be practically useful --- not because the model's algebra changed, but
because the ability to resolve it with this data did. Sparser sampling
(design) can additionally reshape the spectrum, introducing genuinely new
near-degenerate combinations; uniform noise alone only rescales the existing
spectrum, but that alone is enough to push previously-fine directions below
the practical threshold. FIM's projector encodes only the first, permanent
kind of fact (a proven exact symmetry); it has no mechanism for the second
(data-quality-dependent weakness on directions that are structurally fine).

\section{Generalization: how many groups is not the point}
\label{sec:generalization}

A model with two (or more) independent, exactly-isolated symmetry groups,
with a genuinely sharp complement, would be handled by FIM exactly as
well as one group: \code{nullSpaceBasis()}'s threshold classification and
\code{nullSpaceProjector()}'s construction are entirely general in the rank
of the null space, and nothing in the mechanism cares how many dimensions
are exactly zero, only whether they are cleanly separated from everything
else. \emph{The number of exact-symmetry groups is not what breaks the
method; a gapless, sloppy complement is.} This reframes the difference
between PI3K-mini's clean case and the real M-22.0 model precisely: M-22.0
does not fail because it has ``more'' non-identifiability than a toy model
--- it fails because, beyond the one clean PI3K degeneracy, the remaining
$\sim\!91$ parameters show no gap at all (a smooth decline over many orders
of magnitude). A model with three clean PI3K-like degeneracies and an
otherwise sharp, gapped spectrum would likely see FIM succeed just as
cleanly as it did here.

\section{A hybrid regularizer}
\label{sec:hybrid}

Since \code{groupL2()} and \code{constraintL2()} act on orthogonal pieces of
parameter space, they can be summed:
\[
\text{obj\_hybrid} = \text{obj\_data} + \code{groupL2}(P_{\text{null}}, \mu, \lambda_2)
+ \code{constraintL2}(\mu, \sigma_{L2}),\qquad \sigma_{L2}=1/\sqrt{2\lambda_1}.
\]
Equivalently, using $P_{\text{stiff}} := I - P_{\text{null}}$ and two
\code{groupL2} calls with strengths $\lambda_1,\lambda_2$:
\[
\lambda_1 P_{\text{stiff}} + \lambda_2 P_{\text{null}}
= \lambda_1 (I - P_{\text{null}}) + \lambda_2 P_{\text{null}}
= \lambda_1 I + (\lambda_2-\lambda_1) P_{\text{null}}.
\]
This is \emph{always} ``plain isotropic L2 at strength $\lambda_1$, plus a
correction confined to the null direction of size $\lambda_2-\lambda_1$'' ---
the $\lambda_1 I$ term is not conditional on $\lambda_1=\lambda_2$; it falls
directly out of $P_{\text{stiff}}$'s definition as $I-P_{\text{null}}$, for
any $\lambda_1$. Setting $\lambda_1=\lambda_2$ collapses the whole
construction to plain L2 (isotropic penalties are rotation-invariant: equal
weight across \emph{any} complete orthogonal partition of the space
reproduces $\lambda I$, whether that partition is the raw parameter axes or
a null/stiff eigenbasis split). Only $\lambda_1\neq\lambda_2$ produces
something genuinely different from plain L2.

Predicted behavior, following directly from Sections~\ref{sec:clean}--\ref{sec:practical}:
matches FIM's unbiasedness on the one proven-exact direction (strong
$\lambda_2$ dominates there); provides L2's boundedness/crash-resistance
floor everywhere else (weak $\lambda_1$, instead of FIM's exactly-zero
protection there), at the cost of the same small, honest L2-style bias on
those other directions, scaled down by how small $\lambda_1$ is set. This is
not a fix for genuine sloppiness --- the ``everything else'' term still pays
exactly the same data-blind L2 tax plain L2 always has --- it is a more
careful \emph{allocation} of where that tax is paid: never on the one
direction provable to be exact, always (lightly) elsewhere. A three-or-more-
tier version (extra protection specifically for ``moderately sloppy but not
exactly null'' directions, based on their eigenvalue) would re-inherit the
\emph{no-free-lunch} problem already documented in
\file{Technical-Notes\_PI3K-Symmetry-and-Regularization-Diagnostics.pdf}
(\S3--4): there is no principled threshold separating ``weak but real'' from
``weak and worthless'' in a genuinely gapless spectrum. The two-tier version
above is defensible specifically because its boundary sits at an
algebraically \emph{proven} zero, not an arbitrary cutoff.

Not yet implemented/tested as of this note; a natural next step
(\code{obj\_hybrid} as a fourth objective in the fitting script, rerun across
the same $\sigma$ sweep).

\section{Implications for the real HGF model, and chapter strategy}
\label{sec:real-model}

The existing M-22.0 analysis (\file{Chapter-Outline\_FIM-vs-L2-regularization.pdf})
already shows exactly the signature this toy model predicts: one clean,
reproducible, mechanistically-explicable exact degeneracy (PI3K) found
cleanly by the same method, embedded in an otherwise gapless, $\sim\!90/96$-
direction sloppy spectrum. The toy model does not predict a \emph{different}
outcome for the real model; it explains the outcome the real model already
gave, and demonstrates the same transition (clean win $\to$ failure)
happening to one direction at a time, on demand, in a system small enough to
eigendecompose and inspect by hand.

A separate, larger, more biologically comprehensive version of the HGF model
was also discussed as a candidate addition to the chapter (the version
originally requested by the biology collaborators, which needed heavy
hand-tuned priors, converged poorly, and crashed frequently before being cut
down to M-22.0). Two considerations argue against expecting FIM to rescue it,
and against attempting a full three-way multistart comparison on it without
first addressing its numerical robustness separately:

\begin{enumerate}
\item ``Needs heavy priors, converges badly, crashes'' is the signature of
  \emph{many} weakly-constrained directions (the sloppy regime this note
  shows FIM cannot help with), not one isolated exact symmetry.
\item Crashes specifically implicate directions left \emph{completely}
  unconstrained during optimization --- exactly where L2's indiscriminate,
  every-coordinate floor is structurally better suited than FIM's
  single/few-direction-targeted one. Diagnosing which of possibly several
  overlapping near-symmetries a larger model's Hessian encodes, reliably
  enough to trust a projector built from it, is also a substantially harder
  version of the \code{tol}-sensitivity problem already documented for
  M-22.0.
\end{enumerate}

The more defensible use of that model in the chapter is as a motivating,
real-stakes case reported qualitatively (why a smarter regularizer was
wanted in the first place), rather than as a second full quantitative
comparison --- unless its numerical stability is addressed as separate,
scoped work first.

\section{Artifact inventory}

\begin{center}
\begin{tabular}{ll}
\toprule
File & Purpose \\
\midrule
\file{scripts/example\_toy\_pi3kmini\_fim\_fit.R} & model, one multistart per objective, saves \code{.rds} \\
\file{scripts/example\_toy\_pi3kmini\_fim\_plots.R} & waterfall, group scatter, fit-vs-data, \code{plotPars} panels, scatter clouds \\
\file{scripts/example\_toy\_pi3kmini\_fim\_eigenvalue\_report.R} & CSV + plot of initial eigenvalue spectrum vs.\ $\sigma$ \\
\file{scripts/plots/pi3kmini\_fim\_fit/pi3kmini\_results*.rds} & saved parframes/best fits/Hessians, one per $\sigma$/data setting \\
\file{scripts/plots/pi3kmini\_fim\_fit/pi3kmini\_eigenvalue\_report.csv} & eigenvalues at initial point, all ranks, all $\sigma$ \\
\file{2026\_PhdTex/Plots\_fim/pi3kmini\_cartoon.tex/.pdf} & model topology diagram \\
\bottomrule
\end{tabular}
\end{center}

\section{Open items}

\begin{itemize}
\item Implement and test the hybrid regularizer (Section~\ref{sec:hybrid}) across the $\sigma$ sweep.
\item Optional: extend PI3K-mini to two independent symmetry groups to confirm Section~\ref{sec:generalization} directly (not yet run).
\item Optional: a sparser/noisier variant of \code{ab\_fit.R} to confirm Section~\ref{sec:practical}'s AB-model prediction directly (not yet run).
\item Decide, with the above considerations in mind, whether/how to scope any use of the larger biological model in the chapter.
\end{itemize}

\end{document}
